\documentclass[11pt,a4paper]{article}
\usepackage[T1]{fontenc}
\usepackage[utf8]{inputenc}
\usepackage[french]{babel}
\usepackage{lmodern}
\usepackage[margin=2.3cm]{geometry}
\usepackage{amsmath,amssymb}
\usepackage{booktabs,longtable,array}
\usepackage{xcolor,tcolorbox,enumitem,fancyhdr,titlesec}
\usepackage[hidelinks]{hyperref}
\usepackage{siunitx}
\sisetup{group-separator={\,},group-minimum-digits=4,detect-all}

\definecolor{bleufonce}{HTML}{1F3864}
\definecolor{bleuclair}{HTML}{DCE6F1}
\definecolor{grisdoux}{HTML}{F4F4F4}
\definecolor{rougesobre}{HTML}{8B2222}
\titleformat{\section}{\Large\bfseries\color{bleufonce}}{\thesection}{0.7em}{}
\titleformat{\subsection}{\large\bfseries\color{bleufonce}}{\thesubsection}{0.7em}{}
\pagestyle{fancy}
\fancyhf{}
\lhead{\small\color{bleufonce}Traitement des salaires CNPS}
\rhead{\small\color{bleufonce}\thepage}
\renewcommand{\headrulewidth}{0.4pt}
\newtcolorbox{encadre}[1]{colback=bleuclair,colframe=bleufonce,title=#1,
  fonttitle=\bfseries,boxrule=0.6pt,arc=2pt,left=6pt,right=6pt}
\newtcolorbox{alerte}[1]{colback=white,colframe=rougesobre,title=#1,
  fonttitle=\bfseries,boxrule=0.6pt,arc=2pt,left=6pt,right=6pt}
\newtcolorbox{retenir}{colback=grisdoux,colframe=black!30,
  boxrule=0.5pt,arc=2pt,left=6pt,right=6pt}

\begin{document}
\begin{titlepage}
\centering
\vspace*{2.2cm}
{\color{bleufonce}\rule{\textwidth}{1.2pt}}\\[1.0cm]
{\Huge\bfseries Note méthodologique\\[0.3cm]Traitement des déclarations salariales CNPS}\\[0.8cm]
{\Large Pipeline Python v2 corrigé}\\[0.6cm]
{\color{bleufonce}\rule{\textwidth}{1.2pt}}\\[1.5cm]
{\large Unité : salarié--employeur--mois}\\[0.3cm]
{\large Méthode : pondération IPW à deux étages}\\[0.3cm]
{\large Statut : estimations ponctuelles, inférence F.1 en attente}\\[2cm]
\begin{encadre}{Avertissement de diffusion}
Cette version ne revendique aucun intervalle de confiance. Les valeurs réelles
n'ont pas été recalculées pendant l'audit hors VPN. Toute diffusion doit être
rattachée à une session MinIO, à une période, à une empreinte de configuration
et aux contrôles décrits dans cette note.
\end{encadre}
\vfill
{\large 1er août 2026}
\end{titlepage}

\tableofcontents
\newpage

\section{Objet et portée}

Le pipeline transforme les déclarations mensuelles de la CNPS en indicateurs de
distribution salariale. L'unité d'analyse est le couple
\textbf{salarié--employeur--mois}. Une personne employée chez deux employeurs le
même mois forme deux unités légitimes. Dans un tableau cumulé, une personne
présente plusieurs mois contribue plusieurs fois : le tableau décrit des mois
d'emploi, non un stock de personnes uniques.

\begin{encadre}{Estimand}
La cible est la distribution des salaires des unités salarié--employeur--mois
représentables par les employeurs présents dans les fichiers sources et entrant
dans le champ à risque du mois considéré.
\end{encadre}

Sont hors champ : les employeurs absents de tout le panel, les salariés dont
aucune ligne n'existe dans aucune source, et les rémunérations non représentées
par une ligne salarié exploitable. Les résultats ne sont donc pas une mesure
exhaustive de l'emploi ou de la masse salariale de toute l'économie ivoirienne.

\section{Préparation des données}

\subsection{Types, dates et montants}

Le parseur reconnaît les formats numériques français et internationaux. Le
taux d'échec de conversion est journalisé et bloque l'étape au-delà du seuil
configuré. Les âges et anciennetés sont calculés au dernier jour du mois de
déclaration. Ils ne dépendent jamais de la date d'exécution.

\subsection{Déduplication}

Les doublons stricts sont retirés. Pour une clé complète
(individu, employeur, période), la ligne de salaire le plus élevé est conservée.
Les cumuls d'emplois auprès d'employeurs différents sont préservés. Une clé
incomplète n'est jamais dédupliquée : elle est conservée avec un indicateur
explicite de clé de déduplication incomplète.

\subsection{Périodicité salariale}

La variable \texttt{SALAIRE\_BRUT\_ESTIME\_AU\_MOIS} applique les conventions
suivantes : montant mensuel inchangé, taux journalier multiplié par 22,4, taux
horaire multiplié par 179,2. Une périodicité inconnue est traitée comme
mensuelle par défaut; l'hypothèse peut être remplacée par « daily » dans une
analyse de sensibilité.

Le seuil salarial est appliqué dans une unité cohérente avec cette périodicité.
Le volume de périodicités inconnues et le nombre de lignes qu'elles font exclure
sont journalisés. La colonne suffixée \texttt{\_W} est winsorisée après la
conversion. La colonne sans suffixe reste non winsorisée.

\section{Panel entreprise et champ à risque}

\subsection{Réponse de l'employeur}

\[
D_{jt}=\mathbf{1}\{\text{au moins un salaire positif est déclaré par }j
\text{ au mois }t\}.
\]

Un état $D_{jt}=1$ avec salaire moyen non positif, manquant ou non fini est
considéré impossible et provoque un échec.

\subsection{Bornes et absence de registre de cessation}

Le panel commence au mois d'immatriculation, ou à la première apparition si la
date manque. La troncature gauche et le début imputé sont marqués. Aucune
cessation n'est déduite de la dernière déclaration : faute de registre de
radiation, toutes les entreprises sont prolongées jusqu'à la fin commune du
panel et leur fin est dite censurée.

\subsection{Fenêtre glissante}

Pour le mois $t$, une entreprise est dans le champ si elle a déclaré au moins
une fois pendant les $K$ mois \textbf{strictement précédents}. La valeur par
défaut est $K=12$. Le mois courant n'entre pas dans le champ, afin de ne pas
faire dépendre l'univers de la réponse que le modèle doit expliquer.

Les mois d'amorce avec moins de $K$ mois disponibles sont marqués comme fenêtre
extensible. Une diffusion doit comparer $K\in\{6,12,24,\infty\}$.

\subsection{Information disponible à date}

Secteur, commune et taille sont propagés uniquement depuis la dernière valeur
antérieure connue. L'âge de l'entreprise est recalculé à chaque mois. Les
indicateurs « jamais observé auparavant » distinguent une information absente
d'une modalité économique réelle. Aucune valeur future n'est rétropropagée.

\section{Modèles de réponse}

\subsection{Premier étage}

Une régression logistique pénalisée L2 estime
\[
p_{jt}=\Pr(D_{jt}=1\mid X_{jt}).
\]
Les covariables incluent les attributs disponibles à date, la réponse du mois
précédent, le taux de réponse passé et les indicateurs d'amorce.

\subsection{Second étage}

Pour une ligne salarié existante chez un employeur déclarant,
\[
S_{ijt}=\mathbf{1}\{Y_{ijt}>0\},\qquad
q_{ijt}=\Pr(S_{ijt}=1\mid D_{jt}=1,X_{ijt},Z_{jt}).
\]

L'historique est calculé par couple individu--employeur. Un lag n'est qualifié
de mois précédent que pour un écart calendaire égal à un. La complétude de
l'entreprise vient du panel entreprise, est décalée d'un mois civil, puis
attribuée identiquement à toutes les lignes du même couple entreprise--mois.

\begin{alerte}{Limite du second étage}
Le modèle $q$ traite un salaire manquant sur une ligne existante. Il ne peut pas
corriger le salarié dont la ligne entière est absente, faute de registre
exhaustif salarié--employeur--mois.
\end{alerte}

\subsection{Diagnostics}

Les prédictions sont hors échantillon, avec plis groupés par employeur. Sont
contrôlés : classe cible, finitude, recouvrement des propensions,
calibration-in-the-large, pente de calibration, score de Brier, équilibre des
covariables, clipping et trimming. Une strate structurellement jamais
répondante est une violation de positivité.

L'AUC reste descriptive. Une AUC de 0,5 peut être correcte sous MCAR; une AUC
élevée ne garantit pas le recouvrement. Les modèles sont sauvegardés en résumés
JSON non exécutables, jamais en pickle distant.

\section{Pondération et statistiques}

\subsection{Poids final}

\[
R_{ijt}=D_{jt}S_{ijt},\qquad
w^{\mathrm{brut}}_{ijt}=\frac{R_{ijt}}{\widehat p_{jt}\widehat q_{ijt}}.
\]

Le produit positif est tronqué aux quantiles configurés. Les non-répondants et
les lignes hors champ portent un poids nul. Le poids n'est pas stabilisé : il
n'a pas de probabilité marginale au numérateur. Il n'est pas normalisé à
moyenne un. Aucun « effectif pondéré » n'est publié en l'absence de calibration
sur un total de population.

\subsection{Estimateurs}

La moyenne est l'estimateur ratio de Hájek :
\[
\widehat\mu=\frac{\sum_i w_iY_i}{\sum_i w_i}.
\]

La variance diffusée est une dispersion descriptive pondérée avec correction
de Kish; ce n'est pas la variance de $\widehat\mu$. Les quantiles utilisent la
fonction de répartition centrée
\[
F_i=\frac{\sum_{k\le i}w_k-w_i/2}{\sum_kw_k}.
\]
Le Gini est pondéré. Les minimum et maximum sont des extrêmes observés.

La moyenne et la variance descriptive utilisent le salaire winsorisé. Les
quantiles, le Gini et les extrêmes utilisent le salaire non winsorisé. Des
séries mensuelles sont publiées pour le national, le secteur, le sexe et la
taille réduite; les ventilations fines sont libellées « cumul sur la période ».

\section{Inférence et statut des intervalles}

Les déclarations sont traitées comme une population finie affectée par la
non-réponse. Un bootstrap naïf d'employeurs ajouterait une variance de tirage
qui ne correspond pas à ce cadre. L'incertitude vient notamment de
l'estimation conjointe de $p$ et $q$, de leur covariance au niveau employeur et
des non-linéarités des estimateurs.

Cette variance n'est pas encore spécifiée de manière complète. La configuration
impose donc \texttt{point\_only}; chaque sortie porte
\texttt{POINT\_ONLY\_F1\_PENDING}; toute colonne d'intervalle ou d'erreur-type
est rejetée. L'imputation multiple, Rubin, l'AIPW et la double robustesse sont
hors du chemin de publication.

\section{Secret statistique}

Une cellule entière est supprimée si elle contient moins de 30 individus
distincts contributeurs, moins de 3 employeurs distincts, ou si un employeur
représente plus de 85\% de la masse salariale observée non pondérée. Les
identifiants nuls ne comptent pas. Un contributeur a un salaire positif fini et
un poids final positif.

Si une seule cellule est supprimée dans une marge additive, la plus petite
cellule publiée de cette même marge est supprimée en secondaire. Pour les
tableaux mensuels croisés, la règle est appliquée mois par mois. Ces seuils
restent subordonnés à toute doctrine officielle ANStat/CNPS plus stricte.

\section{Validation, export et filiation}

L'export n'est autorisé que si la validation des données, des modèles, des poids
et des cellules passe sans erreur. Le Gini conserve trois décimales dans Excel;
les valeurs nulles ou non finies sont remplacées par un tiret.

Chaque session reçoit un UUID. Son manifeste contient l'empreinte SHA-256 de la
configuration sans secrets, le commit et l'état sale Git, les versions des
dépendances, les durées et sorties déclarées. Les objets canoniques MinIO sont
encore réécrits en place : la réplication historique complète exige leur
versionnement immuable par session.

\section{Contrôles avant diffusion}

\begin{longtable}{p{0.29\textwidth}p{0.64\textwidth}}
\toprule
\textbf{Domaine} & \textbf{Contrôle obligatoire} \\
\midrule
Données & Schéma, parse numérique, clés incomplètes, doublons, périodicité,
volumes avant/après filtre. \\
Temps & Mois continus, âges à date, aucune covariable future, sensibilité $K$. \\
Réponse & Taux $D$ et $S$, support $p/q$, calibration, Brier, équilibre,
clipping et trimming. \\
Poids & Identité $R/(pq)$, poids nuls hors réponse, distribution et extrêmes. \\
Publication & Variables par statistique, séries mensuelles, secret primaire et
secondaire, aucun intervalle F.1. \\
Comparabilité & Pont CIAP versionné, couverture CNPS, rapprochement avec la
rémunération des salariés des comptes nationaux à concepts constants. \\
\bottomrule
\end{longtable}

\section{Limites ouvertes}

\begin{enumerate}[leftmargin=*]
\item variance F.1 et validation de couverture sur processus générateurs;
\item absence de registre de cessation;
\item absence de panel salarié exhaustif pour les lignes entièrement omises;
\item seuils de confidentialité à faire valider institutionnellement;
\item objets MinIO canoniques non immuables;
\item table de passage CIAP--NAPCN à obtenir et versionner officiellement.
\end{enumerate}

\begin{retenir}
Les tests synthétiques valident la logique hors ligne. Ils ne remplacent pas le
recalcul des valeurs réelles, des sensibilités et des rapprochements externes
une fois l'accès VPN/MinIO rétabli.
\end{retenir}

\section*{Références}
\addcontentsline{toc}{section}{Références}
Hájek, J. (1971); Kish, L. (1965); Lerman, R. I. et Yitzhaki, S. (1989);
Lumley, T. (2010); Wooldridge, J. M. (2007).

\end{document}
